% Ad Atticum 10.14 --- headnote
% ad-atticum-10-14, 8 May 49 BC, Cumae

\textit{Cicero to Atticus, written from the Cuman villa
on the eighth day before the Ides of May 49 BC --- 8 May
(the manuscript dateline: \textit{Scr.\ in Cumano viii
Id.\ Mai.\ a.\ 705 (49)}). Servius Sulpicius Rufus,
expected by the Nones (10.9.3) and arrived on the day,
called on Cicero the next morning. The interview is a
study in disappointment. Cicero opens with one of his
sharpest lines on protracted fear: it is a greater evil
to dread a thing so long than to suffer it. Servius
proves the case. He arrives shaken; he fears nothing
that is not worth fearing --- Caesar angry with him,
Pompey no friend, the victory of either side a horror
because of the cruelty of the one and the audacity of
the other, and both sides driven by the same fiscal
desperation, which can be satisfied only out of private
fortunes. He weeps so much that Cicero marvels his tears
have not dried up. The eye-trouble (\textit{lippitudo})
which keeps Cicero from writing in his own hand --- a
copyist takes this letter --- is, by contrast, dry; only
the wakeful nights make it worse.}

\textit{Section 2 asks Atticus for the only kind of
consolation that is now any use: not the philosophical
sort (\textit{id quidem domi est}, ``that resource is at
home'') but news --- of Spain, of Marseilles, of the
two legions Servius has reports on. Section 3 returns to
Servius: their conversation was put off to the next day,
but Servius, slow to set out, said he would rather stay
in his own little bed whatever should come. He clings
firmly to one thing only --- if the condemned exiles
are restored, he will go into exile himself --- and is
deaf to Cicero's arguments that this is already a
foregone conclusion. Cicero concludes that Servius is
to be kept out of his plan rather than drawn into it,
and notes that he will think on Atticus's warning about
Marcus Caelius Rufus.}
